% Source: https://artofproblemsolving.com/community/c3h3590112_varied_algebra
% Stable identity: aops-topic-3590112-q11
After $r$ complete rounds, provided they occur, player $1$ has $k+n-1+r(n-1)$ coins and player $i\ge2$ has $k+n-i-r$. In round $r$ (starting at $r=0$), player $i$ is required to give $rn+i$. After receiving the preceding transfer, every player about to act has $k+n-1+r(n-1)$ coins. Thus that transfer succeeds exactly when $r\le k+n-1-i$. All of the first $k$ rounds succeed. In round $r=k$, players $1,\ldots,n-1$ succeed and player $n$ fails. The final holdings are\[n-2,n-3,\ldots,1,0,n(k+1)-1.\]Two consecutive positive integers cannot be in ratio $5:1$, and a pair containing zero is disallowed. Hence the fivefold pair must be the last and first players, with $n(k+1)-1=5(n-2)$. This is $n(4-k)=9$. With $n\ge2,k\ge1$, the possibilities are $(n,k)=(3,1),(9,3)$. Their initial totals $nk+\dfrac{n(n-1)}{2}$ are $\boxed{6\text{ and }63}$.
